Changing only the training objective of a standard MS lesion segmentation
model halved the number of completely missed lesions on MSLesSeg and
improved the detection of lesions of at least 3~mm on two datasets, without
loss of voxel-level agreement; the gain was reproduced on 47 further
patients and held against four alternative objectives and four alternative
architectures. The ablation attributes this gain to the lesion-level MIL
term. Its cost, additional small false-positive detections, was largely
removed by a 5-voxel size filter on MSLesSeg but was considerably larger on
3D-MR-MS, where a 20-voxel filter was needed.

\paragraph{Clinical relevance of missed small lesions.}
Because the 2024 McDonald criteria rest on counting lesions by location and
over time, and new lesions on follow-up are the main radiological sign of
disease activity (\cref{sec:intro}), the clinically decisive error is often
a single missed lesion rather than an inaccurate lesion volume. This is
also the error that overlap metrics conceal: on one test scan, CATMIL
recovered about 11 lesions that the standard model missed while the
scan-level Dice score did not change (\cref{sec:main_results}). A model
selected and reported by its Dice score can therefore appear accurate while
it misses the lesions that would change a report. For lesions of at least
3~mm, CATMIL raised recall over the standard objective by 6.5--7.4
percentage points, and by about 10 points for the smaller lesions in this
group (\cref{tab:clinical_size}). These lesions are also the ones in which
the central vein sign can be assessed, and one of the recovered lesions in
\cref{fig:qual_lesion_gallery} is juxtacortical, a location that counts
towards dissemination in space. Detecting them therefore affects the
information available for diagnosis, not only the total lesion count.

\paragraph{Use in tandem with a human reader.}
For an assistive tool, a missed lesion is worse than a false positive,
because the reader can dismiss a false positive but receives no prompt for
an unmarked lesion (\cref{sec:intro}). The trade-off produced by CATMIL is
therefore in the favorable direction, provided that the number of false
positives remains manageable. The false positives of CATMIL were small,
lay at a median distance of 1--2~mm from true lesions, and added little
false-positive volume; some of them, such as the focal hyperintensity in
\cref{fig:qual_lesion_gallery}, may be small lesions that the reference
annotation does not contain. Because the minimum component size can be
chosen after training, one trained model can serve different uses: a small
filter for follow-up reading, where sensitivity matters most and every
detection is confirmed by the radiologist, and a larger filter where a
cleaner lesion map is needed, for example for lesion volume measurement.
At the lesion-wise precision of the standard nnU-Net objective, a filtered
CATMIL model still detected more small lesions and had a higher lesion-wise
F1 (\cref{tab:matched_precision}).

\paragraph{Why current models miss small lesions.}
The baselines produced no response at all for most of the small lesions
they missed, so these lesions cannot be recovered by lowering the decision
threshold; the omission has to be addressed during training. Losses that
penalize missed voxels more heavily (Tversky-based) or balance lesions
within an overlap term (CC-DiceCE~\cite{bouteille_learning_2026}) recovered
part of these lesions but remained below CATMIL, and the ablation
(\cref{appendix:catmil-components}) attributes the gain to the lesion-level
detection term: the MIL term alone reached the small-lesion recall of
CATMIL, whereas the CAT term alone, which balances lesions within an
overlap term as CC-DiceCE does, performed like CC-DiceCE. Reweighting
lesions within an overlap loss is therefore not sufficient; missed lesions
have to be penalized directly. The CAT term contributes to how the
detected lesions are delineated: combined with the MIL term, it gave the
best DSC and HD95 of the ablation variants, although within the variation
between models. Changing the network architecture without changing the
objective did not recover them (\cref{sec:arch_vs_objective}). Together
with reports that the errors of current MS segmentation tools are
concentrated in small lesions~\cite{dereskewicz_flames_2025}, this suggests
that such errors are better addressed through the training objective than
through larger networks; because CATMIL changes only the objective, it can
be adopted in existing nnU-Net-based pipelines without changes to
inference.

\paragraph{Limitations.}
Several limitations should be noted. First, this study measures detection
against a reference segmentation, not the effect on radiologists: whether
CATMIL improves the accuracy, agreement, or reading time of human readers
was not tested and requires a reader study. Second, each scan was segmented
independently; the detection of new lesions between serial scans, which is
central to monitoring, was not evaluated directly, and lesion location was
not analyzed, so the effect on the fulfilment of dissemination in space
remains to be shown. Both datasets also cover only the brain, not the
spinal cord or the optic nerve. Third, the evaluation used two public,
retrospective datasets with six test patients each, so the analysis rests
on patient-level confidence intervals rather than significance tests
(\cref{sec:stats}); on 3D-MR-MS, one patient contributes 64\% of the
reference lesions. The 47 MSLesSeg
cross-validation patients reproduce the main findings
(\cref{sec:catmil_oof}), but the corresponding analysis was not performed
for 3D-MR-MS, and validation on a clinical cohort from routine practice is
needed. The ablation of the two CATMIL terms was likewise performed on
MSLesSeg only, so whether the CAT term adds value on other data is
untested. Fourth, the reference segmentations are themselves imperfect: some
detections counted as false positives may be unannotated lesions, and some
missed lesions may be annotations that the image barely supports. Fifth,
the false-positive cost of CATMIL was larger on 3D-MR-MS, where a larger
size filter was needed, and the filter size must be chosen for each
setting rather than being fixed by the method; the operating points at
matched precision were selected on the MSLesSeg test set, so those
MSLesSeg values are optimistic. Finally, lesions of only a few voxels and
lesions with very low contrast remained difficult for every objective: the
method increases the influence of small lesions during training, but it
cannot compensate for evidence that the image does not contain.

\paragraph{Future work.}
Two analyses on existing data come first: the cross-validation analysis
of \cref{sec:catmil_oof} repeated for the 24 3D-MR-MS training patients,
which would reduce the dependence of the 3D-MR-MS results on the single
patient who dominates its test set, and the ablation of the CAT and MIL
terms repeated on 3D-MR-MS. The next clinical step is a reader study in
which radiologists assess MRI scans with and without CATMIL detections,
measuring lesion detection, agreement between readers, and reading time
against an expert consensus reference.
A longitudinal extension that compares serial scans would address the
detection of new lesions directly, and an analysis by McDonald location
would show whether the recovered lesions change the classification of
dissemination in space. Methodologically, a penalty on predicted
components that overlap no reference lesion, as in
BiCC~\cite{bouteille_bicc_2026}, could reduce the false positives within
the training objective rather than by post-processing. The objective may
also be useful for other conditions in which small focal lesions are
clinically important, such as brain metastases.
